% Generated by render.py, do not edit
\documentclass[aspectratio=169,11pt]{beamer}
\usetheme[progress]{presentation}% [progress]: chapter progress bar in the footer

\title{Deck title}
\subtitle{One sentence that says what this talk gives you}
\author{Author}
\institute{Team or institution}% optional
\date{Venue \textperiodcentered{} Year}

\begin{document}

\begin{frame}[plain]
  \titlepage
\end{frame}

\begin{frame}{Agenda}{The question this talk answers in one line}
  \tableofcontents
  \note{What you say to frame the talk: the question it answers and why the chapters come in this order.}
\end{frame}

% Every \section opens a chapter slide, its subsections become that chapter's agenda
\section{Text}
\subsection{Points}

% Bullets: cause and effect points, one message per slide, one grammar per list
\begin{frame}[c]{Claim as a sentence}{Italic context in one line}
  \begin{itemize}
    \item Condition \(\to\) consequence
    \item Second condition \(\to\) second consequence
    \item \hl{Highlight} at most one phrase per slide
  \end{itemize}
  \note{What you say, not what the slide shows: the reasoning behind the points and one example.}
\end{frame}


% Cards: two to four parallel ideas in equal-height cards, each with a serif heading and one or two lines of text
\begin{frame}[c]{Claim that groups parallel ideas}{Use two to four cards, one heading each}
  \begin{cardrow}{3}
    \begin{cardtitled}{First idea}Short text that supports the heading.\end{cardtitled}
    \begin{cardtitled}{Second idea}Short text that supports the heading.\end{cardtitled}
    \begin{cardtitled}{Third idea}Short text that supports the heading.\end{cardtitled}
  \end{cardrow}
  \note{What you say about the ideas: what ties them together and which one matters most.}
\end{frame}

\subsection{Theses}

% Statement: a thesis or a quote on its own slide, the source is optional
\begin{frame}[plain,c]
  \statement[Source or chapter of the thesis]{Thesis as one sentence\\on at most two lines.}
  \note{What you say to land the thesis: why it holds and what the audience should do with it.}
\end{frame}

\section{Numbers}
\subsection{Key figures}

% Hero: one big number carries the slide, one line says what it measures
\begin{frame}[c]{Claim the number proves}{Context of the number in one line}
  \hero{18\texttimes}{what the number measures, against which reference}
  \note{What you say about the number: how it was measured and why it holds beyond this one case.}
\end{frame}


% Key figures: two or three numbers carry the result, the optional [fraction] draws a bar against the full track
\begin{frame}[c]{Claim the numbers prove}{Two or three figures per slide}
  \begin{keyfigures}{3}
    \keyfigure[0.49]{49\%}{what the first number measures}
    \keyfigure[0.78]{78\%}{what the second number measures}
    \keyfigure[0.92]{92\%}{what the third number measures}
  \end{keyfigures}
  \vspace{1cm}
  \subtle{Full track: the reference all three numbers are measured against.}
  \note{What you say about the numbers: how they were measured and how close each gets to the reference.}
\end{frame}

\subsection{Equations}

% Key equation: one equation with every term named, the term in focus highlighted
\begin{frame}[c]{Claim the equation carries}{The highlighted term is the one this talk changes}
  \begin{keyequation}{\mathcal{L} = a + \beta\,\eqhl{b}}
    \term{\mathcal{L}}{quantity}\term{a}{first term}\term{\beta}{weight}\term{b}{term in focus}
  \end{keyequation}
  \note{What you say about the equation: the intuition behind it and why the highlighted term is the one that changes.}
\end{frame}

\section{Tables}
\subsection{Results}

% Data table: options as rows sorted best first, metrics as columns, \bestrow tints the winner
\begin{frame}[c]{Claim the table supports}{Rows are the options, columns the metrics, lower is better}
  \begin{datatable}{X r r r}
    & \hd{Metric 1} & \hd{Metric 2} & \hd{Metric 3}\\
    \midrule
    \bestrow Option C & 2.52 & 7.27 & 1.17\\
    Option B & 3.23 & 7.87 & 1.27\\
    Option A & 6.11 & 9.82 & 5.93\\
  \end{datatable}
  \vspace{\gcgap}
  \muted{One sentence reading the table: how far the winner leads and on which metric.}
  \note{What you say about the table: which row to look at, which direction is better and whether the gap matters.}
\end{frame}

\subsection{Decisions}

% Decision: options against criteria rated with \harvey{0} to \harvey{4}, the last row says when to pick each
\begin{frame}[c]{Claim about the choice}{Filled quarters replace adjectives, the last row is the decision}
  \begin{datatable}{X c c c}
    & \hd{Option A} & \hd{Option B} & \hd{Option C}\\
    \midrule
    First criterion & \harvey{4} & \harvey{2} & \harvey{0}\\
    Second criterion & \harvey{3} & \harvey{4} & \harvey{1}\\
    Third criterion & \harvey{0} & \harvey{2} & \harvey{4}\\
    \midrule
    \muted{Pick it for} & \best{use case A} & use case B & use case C\\
  \end{datatable}
  \note{What you say about the trade-off: which criterion decides and when you would pick another option.}
\end{frame}

\section{Figures}
\subsection{Single figures}

% Full bleed: one image covers the slide, a short claim on a paper band stays legible on any image
\begin{frame}[plain]
  \fullbleed[Source, place or date of the image]{example-image-16x9}{Short claim the image proves,\\at most two lines}
  \note{What you say over the image: where it comes from and the one detail the audience should look for.}
\end{frame}


% Evidence: the claim in words beside the figure that proves it, the right column takes any figure, table or chart
\begin{frame}[c]{Claim the evidence proves}{Setup of the evidence in one line}
  \begin{columns}[c,onlytextwidth]
    \column{0.34\textwidth}
    \claim{One short claim}{One or two sentences on what to see in the figure.}
    \vspace{\gcgap}
    \subtle{Source or setup in one line.}
    \column{0.62\textwidth}
    \roundimage{example-image-16x9}
  \end{columns}
  \note{What you say about the evidence: what to look at first and why it proves the claim, not a reading of the text.}
\end{frame}


% Annotated figure: numbered markers on the figure, notes beside it, revealed one by one
\begin{frame}[c]{Claim the evidence supports}{Numbers on the figure, notes beside it, revealed one by one}
  \begin{annotated}{\roundimage{example-image-16x9}}
    \annot<1->{0.3}{0.62}{First observation on the figure.}
    \annot<2->{0.52}{0.45}{Second observation on the figure.}
    \annot<3->{0.8}{0.3}{Third observation on the figure.}
  \end{annotated}
  \note{What you say at each marker: why the observation matters and how it backs the claim, not a reading of the notes.}
\end{frame}

\subsection{Charts and sets}

% Chart: one pgfplots chart in the house style, direct labels at the line ends, the series in focus in accent
\begin{frame}[c]{Claim the chart proves}{Unit and direction in one line, lower is better}
  \begin{tikzpicture}
    \begin{axis}[gc chart,xlabel={Unit of the x axis},xtick={1,...,6},ymin=0,ymax=100]
      \addplot coordinates {(1,82) (2,80) (3,84) (4,79) (5,81) (6,80)} node[gc label]{Baseline};
      \addplot[gc focus] coordinates {(1,80) (2,64) (3,52) (4,45) (5,41) (6,39)} node[gc label]{New method};
    \end{axis}
  \end{tikzpicture}
  \note{What you say about the chart: what the axes measure, where the lines part and why the gap matters.}
\end{frame}


% Figure row: equal images side by side, same view, same scale, one caption each
\begin{frame}[c]{Claim the comparison shows}{Same view, same scale, one caption each}
  \begin{figrow}
    \fig{example-image-a}{First caption}
    \fig{example-image-b}{Second caption}
    \fig{example-image-c}{Third caption}
  \end{figrow}
  \note{What you say to compare the images: which difference to spot and what it proves.}
\end{frame}

\section{Process}
\subsection{Steps}

% Steps: a sequence on one numbered rail, one click per step, the current step carries the only accent. [0] is a plain pipeline
\begin{frame}[c]{Claim the sequence makes}{Done steps in ink, the current step in accent, the rest still faint}
  \begin{steps}
    \step{First step}{What happens in this step.}
    \step{Second step}{What happens in this step.}
    \step{Third step}{What happens in this step.}
    \step{Fourth step}{What happens in this step.}
  \end{steps}
  \note{What you say at each step: what the done steps delivered and why the current step comes next.}
\end{frame}

\subsection{Timelines}

% Roadmap: phases as segments, milestones as dots, one mark for today
\begin{frame}[c]{Claim about the plan}{Phases as segments, milestones as dots, one mark for today}
  \begin{roadmap}[4.5]
    \phase{3}{First phase}{Months 1 to 3}
    \phase{4}{Second phase}{Months 4 to 7}
    \phase{3}{Third phase}{Months 8 to 10}
    \phase{2}{Last phase}{Months 11, 12}
    \milestone{3}{First milestone}
    \milestone{7}{Second milestone}
    \milestone{12}{Final}
  \end{roadmap}
  \note{What you say about the plan: where we stand today, the next milestone and what could delay it.}
\end{frame}

\begin{frame}[plain,c]
  \statement[Reference or link]{The one sentence\\to remember.}
  \note{What you say to close: the one sentence again in your own words, then the next step for the audience.}
\end{frame}

\end{document}
